\section{课程安排与核心内容}
\subsection{课程安排与阅读材料}
以下课程事项对应 2026 年 9 月 15 日。

\paragraph{Lab 2 提交要求}
Lab 2 在授课当周的星期五截止。需要完成全部要求的代码，包括主机内存的分配与释放；将工作推送到 GitHub；代码完成之后再提交 Lab 2 quiz。

\paragraph{课程辅助资源}
课程聊天机器人可用于询问设备内存分配、主机到设备的数据复制、GPU \term{内核}{kernel}的调用、共享内存数组声明，以及向量加法代码的逐行解释。复习时也可要求它围绕 Lecture 2 和 Lecture 3 逐题提问、等待回答并检查正确性。阅读任务是教材第 6 章和 CUDA Best Practices Guide 的 \textit{Memory Optimizations}；待办事项是提交 Lab 2。

\subsection{学习目标与知识范围}
\paragraph{学习目标}
理解动态随机存取存储器的组织方式；理解突发模式和多个存储体如何提高数据传输率；理解连接 GPU kernel 性能与 DRAM 组织的合并访存；掌握多种 kernel 优化手段及其相互权衡。

\paragraph{优化机制之间的关系}
每块线程数、每块共享内存量和每线程寄存器数共同影响\term{占用率}{occupancy}；减少控制分歧可提高 SIMD 效率，共享内存分块和寄存器分块则减少内存流量。访存与执行优化包括合并访存、面向内存延迟的占用率调节、线程粗化、循环展开和双缓冲。私有化与 warp 级原语分别用于减少访问冲突和线程间协调开销。

全局内存的标称带宽与程序实际获得的有效数据率之间可能存在差距。并发访问若争用同一资源，会产生串行化与延迟；为计算核心提供足够多的线程，还需要让这些线程以合适的方式请求数据。

\begin{keybox}[title={三种互补的改进方向}]
数据复用减少需要搬运的数据量；合并访存提高必要传输中有效数据的比例；并发请求使内存系统在等待某次访问时仍能处理其他工作。线程粗化、指令调度与双缓冲进一步调整每线程工作量和等待关系。
\end{keybox}
